\section{Discussion: Policy Design}\label{sec:discussion}

The policy lesson is not that set-asides are universally dominated. It is that bidder-exclusion rules should be evaluated against the thickness of the protected pool, the heterogeneity of the good, and the welfare weight the planner places on SME surplus. Law 14{,}133/2021 allows both full SME set-asides and price preferences. The estimates here imply that the choice between them should be contingent on market structure rather than imposed as a uniform default.

\subsection{When does a set-aside become defensible?}

The structural decomposition suggests a simple answer. A full set-aside becomes more defensible when three conditions line up: the SME pool is thin enough that protecting it is a first-order objective, the good is heterogeneous enough that the identity of the induced entrants matters for equilibrium selection, and the planner places a high welfare weight on SME surplus. Those are precisely the conditions under which the policy ranking becomes model-sensitive in the data.

Non-pharma sits at the opposite corner. The baseline SME pool is thicker, the goods are more standardized, and the 10 percent price preference is Pareto-superior under both the main equilibrium-selection reading and the strict-invariance benchmark. That is the market class in which bidder exclusion is hardest to defend. Pharma is closer to the knife-edge: the pool is thinner, heterogeneity is higher, and the policy ranking depends on whether the post-policy entrant mix is treated as selected equilibrium participation or forced to equal the pre-policy primitive. The point is not that pharma is special. The point is that thin, heterogeneous markets are exactly where bidder-exclusion rules should be expected to be more sensitive.

The pharmaceutical bifurcation is the paper's central policy claim, not a hedge against an unwelcome result. The paper does not assert universal preference-rule dominance; it identifies the conditions (thin protected pool, heterogeneous good) under which the standard policy ranking reverses, and it documents that those conditions are met in only one of the two Group-65 sub-classes studied. The pharmaceutical welfare ranking under the main specification depends on interpreting the post-policy SME cost distribution as the result of equilibrium selection rather than as a manifestation of latent unobserved heterogeneity that shifts with the policy; under either reading, the policy ranking we report is the analytically derivable one given the assumed primitive, but the deeper question of which interpretation is empirically correct cannot be settled with a single-reform single-platform dataset. Of the eight (class $\times$ type $\times$ period) cells in the structural sample, the policy reversal between V0 and V3 obtains in exactly the two pharma cells under strict invariance and in zero cells under the main specification; non-pharma is robust to both readings. This is a sharp policy partition, not a fragility that the paper papers over: the conditional reading is the contribution, not a qualifier on a universal contribution.

\subsection{The welfare weight implicit in the current regime}

Following the generalized social-marginal-welfare-weight framework of \citet{saezstantcheva2016}, consider a social planner aggregating producer and consumer surplus with weight $w^{\text{SME}}$ on SME producer surplus. The planner is indifferent between the full set-aside (V0) and the 10 percent price preference (V3) when
\begin{equation}\label{eq:welfare_weight}
w^{\text{SME}}_\star
\;=\; \frac{\mathrm{Loss}(V_0) - \mathrm{Loss}(V_3)}
           {\mathrm{SMESurplus}(V_0) - \mathrm{SMESurplus}(V_3)}.
\end{equation}
Under the main specification, the loss differential $\mathrm{Loss}(V_0) - \mathrm{Loss}(V_3)$ is approximately $0.22$ in non-pharmaceuticals and $0.31$ in pharmaceuticals: V0 welfare loss at $\lambda = 0.30$ is $0.223$ and $0.308$ from Table~\ref{tab:v3_welfare} (rows non-pharma $\lambda = 0.30$ and pharma $\lambda = 0.30$), and V3 welfare loss is near zero from the price-preference simulation in Table~\ref{tab:v3_apv} ($\Delta = -0.0036$ and $+0.0015$ in non-pharma and pharma respectively). The SME-surplus differential $\mathrm{SMESurplus}(V_0) - \mathrm{SMESurplus}(V_3)$ is approximately $0.09$ in non-pharmaceuticals and $0.12$ in pharmaceuticals: under V0 the implicit transfer to SME winners is $\Delta_{\text{gov}} - \mathrm{DWL}_{\text{alloc}} = 0.242 - 0.150 = 0.092$ in non-pharma and $0.328 - 0.210 = 0.118$ in pharma (Table~\ref{tab:v3_welfare}, columns $\Delta_{\text{gov}}$ and DWL$_{\text{alloc}}$).\footnote{Under the English-reverse IPV interpretation of Section~\ref{sec:model}, all losers exit at cost and earn zero auction-stage profit; SMESurplus thus reduces to the winner-conditional rectangular transfer, which is what is reported here. Under V3, the SME-surplus transfer is taken as approximately zero given the near-zero $\Delta$ on the open-regime auction.} Substituting in equation~\eqref{eq:welfare_weight} delivers $w^{\text{SME}}_\star \approx 0.223/0.092 \approx 2.42$ in non-pharmaceuticals and $0.308/0.118 \approx 2.61$ in pharmaceuticals. Under the strict-invariance benchmark, the non-pharmaceutical preference ranking is unchanged (the implied weight remains above the utilitarian benchmark and within the upper end of the Brazilian transfer-program range), while the pharmaceutical indifference weight falls to 0.7 (Table~\ref{tab:v3_strict_invariance}). These are not contradictory numbers. They are a map of when bidder exclusion is doing something that a moderate preference rule cannot.

The comparison is informative because it disciplines the distributional argument in favor of set-asides. A planner willing to pay the non-pharma fiscal cost of bidder exclusion would have to value one real of SME surplus at more than twice a real of general welfare, even in a thick market where a preference rule achieves nearly the same allocative objective. In pharma, the required weight depends on how strongly one believes the policy selects a different entrant mix. The right reading is therefore conditional: bidder exclusion is hardest to justify in thick standardized markets, and most plausible only in thin heterogeneous markets with very high welfare weights on protected firms. Figure~\ref{fig:v3_welfare_weight} plots the implied weights against the Brazilian social-policy benchmark.

\begin{figure}[!htbp]
\centering
\includegraphics[width=0.98\textwidth]{../output/figures/fig_v3_welfare_weight.pdf}
\caption[Implicit welfare weight to prefer the set-aside over the price preference.]{\textit{Implicit welfare weight $w^{\text{SME}}_\star$ required to prefer the full set-aside over a 10 percent price preference, by class and specification.} The grey band marks the 1.2--1.5 range typical of Brazilian cash-transfer programs; the dashed vertical line marks the utilitarian benchmark $w = 1$. In non-pharmaceuticals, the implied weight sits above the transfer-program band under the main equilibrium-selection specification ($w^{\text{SME}}_\star = 2.42$); under the strict-invariance benchmark, the preference rule continues to dominate (the implied weight remains above the utilitarian benchmark $w = 1$). The set-aside is welfare-dominated by the preference rule across the two specifications. In pharmaceuticals, the implied weight under the main specification ($w^{\text{SME}}_\star = 2.61$) sits above the band but reverses to 0.7 under strict invariance, below the utilitarian benchmark; the pharmaceutical ranking is therefore conditional on how the post-policy SME pool is modeled.}
\label{fig:v3_welfare_weight}
\end{figure}

\subsection{Implementation considerations}

Two features qualify the policy comparison. First, V3 is computed in a Vickrey-equivalent framework. In a first-price sealed-bid auction with the same 10 percent preference rule, bidder equilibrium behavior would adjust as characterized by \citet{maskinriley2000}: SME bidders bid less aggressively (because the preference reduces the marginal cost of a higher bid) and non-SMEs bid more aggressively (because they face a competitive disadvantage). For a 10 percent preference, the net effect on the second-order statistic from these two adjustments is bounded below 5 percent of the reference price in expectation. The direction of the adjustment preserves the main insight here---keeping non-SMEs active disciplines the price distribution---and the conditional welfare ranking is robust to it in non-pharma. In pharma the close comparison may be sensitive to it; the strict-invariance reversal is therefore a more first-order concern than the FPSB equilibrium adjustment.

Second, a preference rule is procedurally more complex to implement than a set-aside. The buyer must score bids twice and disclose the winner-selection rule clearly. But this is an operational complication, not a conceptual obstacle. Preference systems have been run elsewhere for decades, and Brazilian procurement law already contemplates them. The relevant policy question is therefore not whether the instrument exists, but in which markets it should replace bidder exclusion.

\subsection{Market thickness and goods heterogeneity}

The non-pharma versus pharma comparison is best read as a statement about the interaction between bidder exclusion and goods characteristics. Pharma auctions have both a thinner SME pool and higher within-class heterogeneity. Once non-SMEs are removed from price formation, those two features raise the expected second-order statistic and make the selected entrant mix more important. Non-pharma auctions are closer to the textbook standardized-good environment in which a moderate preference rule can tilt allocation toward SMEs without materially changing price formation. Operationally, the bifurcation in the data here obtains where post-policy SME-per-auction count is below approximately 1.3: non-pharma at 1.87 SMEs per auction lies above this threshold and the V3 $>$ V0 ranking holds across both specifications; pharma at 1.22 SMEs per auction lies below it and the ranking reverses under strict invariance.

The pattern is consistent with the Group-65 split but the paper does not estimate it as an external-validity claim outside the BEC platform. Within structurally similar Brazilian procurement settings, goods with deeper supplier bases and tighter specification standardization are plausible candidates for preference rules rather than exclusion; goods with thinner protected pools, more heterogeneous production technologies, and stronger administrative reasons to value SME surplus are the cases in which the policy comparison becomes more model-sensitive. Comparable extension settings include federal Brazilian procurement on ComprasNet (a larger universe with the same Law 14{,}133 framework), municipal procurement under similar SME defaults, and non-Brazilian SME-set-aside reforms such as the U.S. SBA 8(a) program changes or EU Directive 2014/24 transposition events, where the same conditional-on-thickness mechanism could be tested against alternative institutional baselines. The contribution of the paper is to put that intuition on structural objects in one specific reform: the size of the within-auction component, the magnitude of the participation-margin offset, and the welfare weight required to rationalize the policy under each specification.

Two further structural questions---the IPV restriction in pharmaceutical procurement and the possibility of confounding regulatory events at the March 2018 cutoff---are taken up in Appendices~\ref{app:robustness_ipv} and~\ref{app:cutoff} respectively. Both qualify the pharmaceutical welfare-dominance result rather than the non-pharmaceutical headline; the pharmaceutical claim is conditional on the equilibrium-selection treatment of $F_c^{\text{SME,Post}}$, while the non-pharmaceutical claim is preserved across the IPV-relaxed bounds and the strict-invariance specification.

\subsection{Extensions the paper does not attempt}

Three extensions represent natural follow-on work. (i) Quantity response: the buyer's adjustment to higher prices, which the static model treats as inelastic at observed values. At a typical demand elasticity for essential medical supplies of $\varepsilon \in [0.1, 0.3]$, the quantity adjustment to the simulated 11 percent log-price increase would shave the fiscal cost by $\varepsilon \cdot 0.11$, or roughly 1--3 percentage points of $p_{S_1}$, and would slightly reduce the allocative DWL component. (ii) Dynamic capacity accumulation: SME firms protected by the rule may build production capacity that persists beyond the policy window, an externally validated mechanism in \citet{ferraz2016} and \citet{szerman2023}. (iii) General-equilibrium SME surplus: the partial-equilibrium SMESurplus measure used in the welfare-weight identity (Section~\ref{sec:discussion}) could be extended to a general-equilibrium framework with option value, capacity persistence, and private-sector spillovers. Each is a margin a fully dynamic-general-equilibrium model would price; each would tend to enlarge the SME-side benefit and could shift the welfare-weight comparison. None changes the organizing lesson of this paper: the welfare ranking of bidder-exclusion rules is conditional on market structure.
